\section*{الفصل \ref{ch:g8:puberty} --- البلوغ}

\begin{solution}{exo:g8:puberty:1}
المرحلة التي يصير فيها جسم الطفل جسمَ بالغ --- إذ تنضج أعضاء
التكاثر وتبدأ إنتاج \omterm{def:g8:sexual-reproduction:gametes}{الأمشاج}، ويأخذ الجسم صورته البالغة. وهو اللبّ
البيولوجي للمراهقة.
\end{solution}

\begin{solution}{exo:g8:puberty:2}
الفتيات: ينمو الثديان، ويتسع الحوض، وتأتي أول الحيضات (مع شعر
الجسم). والفتيان: يغلظ الصوت، وتعرض الكتفان، ويبدأ إنتاج النطف (مع
شعر الوجه). والجنسان: دفعة النمو، \omterm{def:g1:the-five-senses:sense}{وجلد} أدهن، وعرق متغير، وطقس
مزاج.
\end{solution}

\begin{solution}{exo:g8:puberty:3}
عند الفتيات، أول الحيضات --- وهي بدء الدورات؛ وعند الفتيان، بدء
الخصيتين إنتاج النطف.
\end{solution}

\begin{solution}{exo:g8:puberty:4}
رسول كيميائي يطلقه في \omterm{prop:g4:journey-of-food:absorb}{الدم} عضو متخصص، ويحمله الدوران إلى كل مكان،
وتطيعه الأعضاء المبنية للاستجابة له.
\end{solution}

\begin{solution}{exo:g8:puberty:5}
تبدأ عند مركز تحكم في قاعدة \omterm{def:g5:brain-in-command:system}{الدماغ}، وهرموناته توقظ أعضاء التكاثر؛
ثم تأمر هرموناتها هي بإعادة بناء الجسم كله --- فتنتهي إذن في كل
مكان يبلغه \omterm{prop:g4:journey-of-food:absorb}{الدم}.
\end{solution}

\begin{solution}{exo:g8:puberty:6}
كلاهما طبيعي: فالجدول الزمني شخصي --- والمبكر والمتأخر كلاهما
سليم، وترتيب التغيرات متغير، والوجهة مشتركة. وتباعدهما سنوات في
الثالثة عشرة لا يكاد يقول شيئًا عن أي منهما.
\end{solution}

\begin{solution}{exo:g8:puberty:7}
لأن الرسول يسافر في \omterm{prop:g4:journey-of-food:absorb}{الدم}، وهو يزور كل عضو على كل حال: فبثّ واحد
يبلغ صندوق الصوت والكتفين \omterm{def:g1:the-five-senses:sense}{والجلد} سواءً، بلا حاجة إلى أسلاك.
والأعضاء التي تستجيب هي المبنية لسماع ذلك الرسول --- فالعنونة
بالمستقبِل لا بالطريق.
\end{solution}

\begin{solution}{exo:g8:puberty:8}
السبب الأول: هرمونات \omterm{def:g8:puberty:puberty}{البلوغ} نفسها تبلغ \omterm{def:g5:brain-in-command:system}{الدماغ}. والسبب الثاني:
\omterm{def:g5:brain-in-command:system}{الدماغ} يعيد توصيل نفسه في الوقت نفسه نحو صورته البالغة. والطمأنة:
إنه طقس --- طبيعي، مفهوم، وزائل.
\end{solution}

\begin{solution}{exo:g8:puberty:9}
مادة البناء قبل كل شيء --- فالدفعة تبني عظمًا وعضلًا --- مع بنّائي
\omterm{def:g3:bones-and-muscles:skeleton}{العظام} في الألبان \omterm{prop:g3:balanced-diet:jobs}{وطاقة} ثابتة من الأسرة النشوية: وهو منطق الحمية
القاسية في \cref{exo:g7:digestion-and-nutrients:11}، مُجرًى إلى
الأمام.
\end{solution}

\begin{solution}{exo:g8:puberty:10}
ورشة البناء أشدّ ما تكون هشاشةً في وسط البناء: فالتبغ والكحول
وسائرها تضرب أعضاءً --- \omterm{def:g5:brain-in-command:system}{والدماغ} أولها --- ما زالت أسلاكها وأنسجتها
تُمدّ. والضرر في بناء منجَز إصلاح؛ والضرر في بناء نصف قائم تغيير
في التصميم.
\end{solution}

\begin{solution}{exo:g8:puberty:11}
الجدير بالثقة: ممرضة المدرسة، وطبيب الأسرة (مدرَّبان، وسرّيان، ولا
يدهشهما شيء)، وهذا الكتاب (مبني ليُقرأ من جديد ويُراجَع). وغير
الجدير: المقطع المنتشر (فهو يبيع، ويتاجر بالشك)، وأسطورة الملعب
(تخمينات معادة بلا عنوان للتصحيحات).
\end{solution}

\begin{solution}{exo:g8:puberty:12}
المدى: رسول واحد يغيّر الصوت \omterm{def:g1:the-five-senses:sense}{والجلد} والهيكل والمزاج دفعةً واحدة ---
\omterm{prop:g4:journey-of-food:absorb}{فالدم} يوصّل إلى كل مكان (والرصد: إعادة البناء المنسَّقة). والتنسيق:
التغيرات تجري على جدول واحد مع أن لا \omterm{def:g5:brain-in-command:system}{عصب} يصل المواضع (والرصد:
اللحية والصوت يصلان معًا). والاختلاف عن \omterm{def:g5:brain-in-command:system}{الأعصاب}: لا أسلاك، ولا
سرعة جزء من الثانية --- فالأوامر تستغرق شهورًا وتعمل في الجسم كله،
حيث تستغرق رسائل \cref{prop:g5:brain-in-command:loop} أجزاءً من
الألف من الثانية وطريقًا واحدًا (والرصد: سنوات \omterm{def:g8:puberty:puberty}{البلوغ} البطيئة في
مقابل خُمس ثانية اللقفة).
\end{solution}

\begin{solution}{pb:g8:puberty:1}
\textbf{1.} الأربعة كلها تُظهر دفعة \omterm{def:g3:stages-of-life:stages}{المراهقة}. أ: الدفعة تكاد تكتمل،
وهي تستوي نحو طول البالغ. وب: الدفعة لم تبدأ بعد --- ما زال على
صعود \omterm{def:g3:stages-of-life:stages}{الطفولة} الثابت. وج: في وسط الدفعة. ود: دفعة مبكرة، جرت من
التاسعة إلى الثانية عشرة، وهي تستوي الآن.

\textbf{2.} ``منحناك بالضبط منحنى طبيعي قبل الدفعة: فالجدول الزمني
شخصي، والتأخر سليم كالتبكير، وصعودك الثابت يقول إن الدفعة قادمة ---
والمنحنيات مثل منحناك تشتدّ ميلًا في غضون سنة أو سنتين، وتصل إلى
أطوال البالغين نفسها التي يصل إليها الجميع.''

\textbf{3.} الاستواء ينهي دفعة \emph{الطول} --- أي النمو في
القامة. أما بقية عمل \omterm{def:g8:puberty:puberty}{البلوغ} --- الامتلاء، والنضج، وترميم \omterm{def:g5:brain-in-command:system}{الدماغ} ---
فيستمر بعده؛ فصاحب المنحنى د مبكر لا منتهٍ.

\textbf{4.} لأن الطول يخلط الميراث بالعمر --- فطفل طويل قبل الدفعة
قد يفوق في القياس طفلًا قصيرًا في وسطها. أما الميل فيُظهر العملية:
فالحدّة تقول دفعة، والاستواء يقول نهايتها، مهما يكن ارتفاع الخط.

\textbf{5.} استيقاظ مركز التحكم في قاعدة \omterm{def:g5:brain-in-command:system}{الدماغ}؛ فأمرت هرموناته
أعضاء التكاثر بالاستيقاظ، وأمرت هرموناتها --- مبثوثةً \omterm{prop:g4:journey-of-food:absorb}{بالدم} ---
صفائح النمو وسائر الجسم.

\textbf{6.} \omterm{prop:g1:healthy-habits:needs}{النوم}: فالنمو وترميم \omterm{def:g5:brain-in-command:system}{الدماغ} يجريان ليلًا --- والليالي
القصيرة تجوّعهما. والغذاء: فالدفعة تُدفع مادةً --- ومادة البناء
وأسر الصحن تموّلها. (والحركة وترك التخريب يكملان قائمة العناية.)

\textbf{7.} \omterm{def:g8:puberty:hormones}{الهرمونات} تبلغ \omterm{def:g5:brain-in-command:system}{دماغ} ج، \omterm{def:g5:brain-in-command:system}{والدماغ} في وسط إعادة توصيله ---
كيمياء مع بناء، وهو طبيعي وزائل. وعمليًّا: الساعة تنجرف متأخرة
والمدرسة لا تنجرف؛ فعلى ج أن يدافع عن ساعاته بهدأة مبكرة خالية من
ضوء الشاشات
(\cref{prop:g5:healthy-choices:screens}).

\textbf{8.} لأن غرفة الاستشارة مدرَّبة وسرّية ولا تدهشها الأسئلة
--- فالأسئلة التي لا يستطيع المنحنى إظهارها لها عنوان صحيح مع ذلك،
وأسطورة الملعب وأركان الشابكة ليست هي.

\textbf{9.} أ: الدفعة تنتهي، طبيعي --- احرس \omterm{prop:g1:healthy-habits:needs}{النوم} في سنوات
الامتحانات. وب: قبل الدفعة، طبيعي --- صبر وصحن. وج: في وسط الدفعة،
طبيعي --- ساعات \omterm{prop:g1:healthy-habits:needs}{نوم} ومادة بناء. ود: دفعة مبكرة مكتملة، طبيعي ---
حركة لسنوات الامتلاء.

\textbf{10.} لأن الطبيعي طريق واسع لا خط: أربعة جداول سليمة،
تتباعد سنوات، وكلها تصل.

\textbf{11.} لا شيء: فالآلة سليمة وجدولها شخصي --- فلا عطل يُصلَح،
ولا رافعة تستحق الشدّ على منحنى سليم؛ والوصفة الصادقة هي الوقت (مع
قائمة العناية ريثما يمضي).

\textbf{12.} ``\omterm{def:g8:puberty:puberty}{البلوغ} إعادة بناء الجسم من طفل إلى بالغ بأمر
\omterm{def:g8:puberty:hormones}{الهرمونات} --- يديره رسل في \omterm{prop:g4:journey-of-food:absorb}{الدم}، على جدول شخصي، ويموّله \omterm{prop:g1:healthy-habits:needs}{النوم}
والغذاء والحركة.''
\end{solution}
